\documentclass[11pt]{article}
\usepackage{fontspec}
\usepackage[french,english]{babel}
\usepackage{amsmath}
\usepackage{amssymb}
\usepackage{amsthm}
\usepackage{mathrsfs}
\usepackage{tikz-cd}
\usepackage{enumitem}

\newcommand{\printedpage}[1]{\marginpar{\footnotesize p.~#1}}
\newcommand{\msp}[1]{\marginpar{\footnotesize [#1]}}
\newcommand{\uCh}{\underline{\mathrm{Ch}}}
\newcommand{\uHom}{\underline{\mathrm{Hom}}}
\newcommand{\uFl}{\underline{\mathrm{Fl}}}
\newcommand{\bDelta}{\boldsymbol{\Delta}}
\newcommand{\Hom}{\mathrm{Hom}}

\begin{document}

\printedpage{14}
% head: A. Grothendieck : Dérivateurs, Ch. V, Édition des Universités de Montpellier II et Paris VII
What precedes justifies our reasoning, which gave by composition of paths well-defined composition maps
\[
\underbrace{\Hom_\Sigma(x,y)}_{=\pi_0\uCh_\Sigma(x,y)} \times \underbrace{\Hom_\Sigma(y,z)}_{=\pi_0\uCh_\Sigma(y,z)} \longrightarrow \underbrace{\Hom_\Sigma(x,z)}_{=\pi_0\underline{\textstyle\mathrm{Ch}}_\Sigma(x,z)}.
\]
% (the third "Ch" under the brace is printed in a larger size than the two others)

\msp{page 22}I say that this composition is associative. This follows from the fact that if we have three paths (of \textit{even} indices!)
\[
x \xrightarrow{c} y \xrightarrow{c'} z \xrightarrow{c''} t,
\]
we have (without having to carry out homotopies \dots) strict associativity
\[
c''(c'c) = (c''c')c.
\]

I say that there are identities. We define
\[
1_x \in \Hom_\Sigma(x,x)
\]
by the path of index 0 (which is even!), i.e. by the diagram reduced to $x$, with no arrow (or if we want to interpret the diagram as a functor $e \longrightarrow X$, there is the implicit identity arrow). Our arguments did not exclude the case $n = 0$ or $m = 0$, and it is tautological that for a path
\[
x \xrightarrow{c} y
\]
we have
\[
c \cdot 1_x = 1_y \cdot c = c.
\]
So we have indeed defined a category, which I am going to call $X_\Sigma$. I am going to define a canonical functor % [sic] Ja -> Je, misprint of the pronoun
\[
\theta : X \longrightarrow X_\Sigma
\]
which will be the identity on objects, we must define it on arrows
\[
\Hom_X(x,y) \longrightarrow \Hom_{X_\Sigma}(x,y) = \pi_0\mathrm{Ch}_\Sigma(x,y)
\]
\texttt{[rather} $\pi_0\uCh_\Sigma(x,y)$\texttt{]}. Unfortunately, an arrow
\[
x \xrightarrow{u} y
\]
\msp{page 23}cannot be interpreted directly as a path in our sense, we have to pay for our evasions! But we are going to associate to such an arrow the path
\[
x \longrightarrow y \xleftarrow{\mathrm{id}} y \xrightarrow{\mathrm{id}} y \xleftarrow{\mathrm{id}} y,
\]
% foot: Version 9 septembre 2001
\printedpage{15}
% head: A. Grothendieck : Dérivateurs, Ch. V, Édition des Universités de Montpellier II et Paris VII
an element of $\uCh_{\Sigma,2}(x,y)$, which defines the desired map. We must check compatibility with identities and under composition. To the arrow $1_x : x \longrightarrow x$ is associated the element of $\Phi_2(x,x)$
\[
x \xrightarrow{\mathrm{id}} x \xleftarrow{\mathrm{id}} x \xrightarrow{\mathrm{id}} x \xleftarrow{\mathrm{id}} x,
\]
we must prove that in $\pi_0(\uCh_{\Sigma,2}(x,x))$ it becomes equal to the element deduced from $1_x \in \Phi_{\Sigma,0}(x,x)$, by the transition functor. But this element is precisely none other than
\[
x \xrightarrow{\mathrm{id}} x \xleftarrow{\mathrm{id}} x \xrightarrow{\mathrm{id}} x \xleftarrow{\mathrm{id}} x,
\]
it is the same (without even having to pass to $\pi_0$)!

For compatibility with composition
\[
x \xrightarrow{u} y \xrightarrow{v} z,
\]
we have on the one hand
\[
x \xrightarrow{vu} z \xleftarrow{\mathrm{id}} z \xrightarrow{\mathrm{id}} z \xleftarrow{\mathrm{id}} z,
\]
on the other hand
\[
x \xrightarrow{u} y \xleftarrow{\mathrm{id}} y \xrightarrow{\mathrm{id}} y \xleftarrow{\mathrm{id}} y \xrightarrow{v} z \xleftarrow{\mathrm{id}} z \xrightarrow{\mathrm{id}} z \xleftarrow{\mathrm{id}} z.
\]
We complete the first path by degeneracy at both ends, to compare it with the second as an element of $\uCh_4(x,z)$,
\[
x \xrightarrow{\mathrm{id}} x \xleftarrow{\mathrm{id}} x \xrightarrow{vu} z \xleftarrow{\mathrm{id}} z \xrightarrow{\mathrm{id}} z \xleftarrow{\mathrm{id}} z \xrightarrow{\mathrm{id}} z \xleftarrow{\mathrm{id}} z.
\]

\msp{page 24}We must prove that the two are homotopic. We thus have two functors
\[
f,g : \underbrace{\overbrace{\uFl(X;x,y)}^{\text{discrete category}} \times \overbrace{\uFl(X;y,z)}^{\text{discrete category}}}_{\mathcal{A}} \underset{g}{\overset{f}{\rightrightarrows}} \underbrace{\uCh_{4,\Sigma}(X;x,z)}_{\mathcal{B}}, \tag*{$(\ast)$}
\]
we want to prove that they are homotopic. For this, we regard them (as a moment ago) as ``transposes'' of the two functors
\[
\begin{tikzcd}[column sep=large]
\bDelta_1 \sqcup \bDelta_1 \simeq \bDelta_2 & \Phi_4 \arrow[l, shift right, "\varphi"'] \arrow[l, shift left, "\psi"] \\
& \varepsilon, \arrow[ul, "\alpha"] \arrow[u, "\alpha'"']
\end{tikzcd}
\]
since the categories $\mathcal{A}$, $\mathcal{B}$ of the two sides of (*) can be interpreted as full subcategories of $\uHom(\bDelta_2, X)$ and $\uHom(\Phi_4, X)$ respectively; namely the full subcategories consisting of the functors whose composite with $\alpha$ resp. $\alpha'$ is given by $x$ and by $z$, and moreover (for $\mathcal{A}$) whose value at the middle object of $\bDelta_2$ is $y$. We can also replace $\mathcal{A}$ by the larger full subcategory $\mathcal{A}'$ of $\uHom(\bDelta_2, X)$, obtained by ignoring this last condition, and it suffices to prove that the two ``bigger'' functors
% foot: Version 9 septembre 2001
\printedpage{16}
% head: A. Grothendieck : Dérivateurs, Ch. V, Édition des Universités de Montpellier II et Paris VII
\[
\mathcal{A}' \underset{g'}{\overset{f'}{\rightrightarrows}} \mathcal{B}'
\]
are homotopic. And for this, it suffices to
\msp{page 25}see again that $\varphi$ and $\psi$ are homotopic in the category of functors $\Phi_4 \longrightarrow \bDelta_2$ which are compatible with the source, and the target. We thus see that our homotopy lemma of a moment ago was not general enough. We would have to define the \textit{generalized interval categories,} of arbitrary length $n$ ($n =$ number of arrows end to end), a \textit{standard} generalized interval category of length $n$ has as its set of objects the interval of integers % [sic] la catégories -> la catégorie, singular article with plural noun
\[
[0,n],
\]
and for every pair of consecutive integers $i-1$, $i$ among the objects ($1 \leq i \leq n$), there is exactly one arrow connecting $i$, $i+1$ \texttt{[rather} $i-1$\texttt{,} $i$\texttt{]}. We assign to the interval $[i,i+1]$ the sign $+$ if $i \longrightarrow i+1$ (arrow in the direct direction), the sign $i$ \texttt{[rather} $-$\texttt{]} if $i \longleftarrow i+1$ (arrow in the retrograde direction) - thus the standard categories are defined, for fixed $n$, by a system of signs $+$ or $-$, which we assign to the $n$ consecutive ``segments'' $A_i := [i-1,i]$ ($1 \leq i \leq n$) of the interval of integers $[0,n]$. There are thus exactly $2^n$ of them.

We expect the following statement to be true:

\textbf{Lemma.} \textit{Let $I$, $I'$ be two generalized intervals, equipped with source objects and target object (so that $I$, $I'$ are} canonically \textit{isomorphic to the corresponding standard generalized intervals). Let} $\uHom^!(I,I')$ % [sic] genéralisés -> généralisés, missing accent on the first e
\msp{page 26}\textit{be the category of functors from $I$ to $I'$, compatible with source and target. If this category is $\neq \emptyset$, it is} 0\textit{-connected} (and, I presume, even contractible?).

I am going to assume provisionally that this is so, at least in the particular case of $\Phi_4$ (generalized interval of length 8) and $\bDelta_2$. This therefore gives a functor
\[
\theta : X \longrightarrow X_\Sigma.
\]
We have won, for the universal property of this functor is practically obvious. But I forgot to check that $\theta$ transforms elements of $\Sigma$ into invertible arrows. So let
\[
x \xrightarrow{u} y
\]
with $u \in \Sigma$, we are going to define a path in the reverse direction
\[
y \xrightarrow{c} x
\]
which will define in $X_\Sigma$ an arrow inverse to $u$. The natural idea is to take the ``path''
\[
y \xleftarrow{u} x
\]
% foot: Version 9 septembre 2001
\printedpage{17}
% head: A. Grothendieck : Dérivateurs, Ch. V, Édition des Universités de Montpellier II et Paris VII
from $y$ to $x$, but this is not a path in the sense of our definitions, we are going to complete it to
\[
c : y \xrightarrow{\mathrm{id}} y \xleftarrow{u} x \xrightarrow{\mathrm{id}} x \xleftarrow{\mathrm{id}} x.
\]
I say that $c$ does the job. We must prove two homotopy relations, of $c \cdot \bar{u}$ and of $\bar{u} \cdot c$ with the identity paths at $x$ \texttt{[resp. at} $y$\texttt{]} of length 8 (index 4, i.e. in $\uCh_{\Sigma,4}(X;x,x)$ resp. $\uCh_{\Sigma,4}(X;y,y)$),
\msp{page 27}where $\bar{u}$ is the path associated to $u$
\[
x \xrightarrow{u} y \xleftarrow{\mathrm{id}} y \xrightarrow{\mathrm{id}} y \xleftarrow{\mathrm{id}} y.
\]

We thus have to look at
\[
x \xrightarrow{u} y \xleftarrow{\mathrm{id}} y \xrightarrow{\mathrm{id}} y \xleftarrow{\mathrm{id}} y \xrightarrow{\mathrm{id}} y \xleftarrow{u} x \xrightarrow{\mathrm{id}} x \xleftarrow{\mathrm{id}} x
\]
and see that it is homotopic to the trivial path
\[
x \xrightarrow{\mathrm{id}} x \xleftarrow{\mathrm{id}} \cdots x \xrightarrow{\mathrm{id}} x \xleftarrow{\mathrm{id}} x
\]
(identities everywhere), and likewise for
\[
y \xrightarrow{\mathrm{id}} y \xleftarrow{u} x \xrightarrow{\mathrm{id}} x \xleftarrow{\mathrm{id}} x \xrightarrow{u} y \xleftarrow{\mathrm{id}} y \xrightarrow{\mathrm{id}} y \xleftarrow{\mathrm{id}} y
\]
and the trivial path at $y$. This time, we again have two functors
\[
\begin{tikzcd}[column sep=large]
\bDelta_1 & \Phi_4 \arrow[l, shift right, "\varphi"'] \arrow[l, shift left, "\psi"] \\
& \varepsilon, \arrow[ul, "\alpha'"] \arrow[u, "\alpha"']
\end{tikzcd}
\]
compatible with source objects and target objects, and we must prove that they are homotopic in $\uHom^!(\Phi_4, \bDelta_1)$. This must therefore follow from the same general lemma. We have won - by a simple interpretation of $X\Sigma^{-1}$, and in particular of the $\Hom_{X\Sigma^{-1}}(x,y)$ and of their composition, in terms of the geometry of path spaces:
\[
\left\{
\begin{aligned}
&\Hom_{X\Sigma^{-1}} \simeq \pi_0(\uCh_\Sigma(X;x,y)) \\ % [sic] Hom_{XΣ^{-1}} -> Hom_{XΣ^{-1}}(x,y), arguments (x,y) missing
&\text{Composition of homomorphisms in } X\Sigma^{-1} \text{ corresponds to} \\
&\text{the composition of } \Sigma\text{-paths in } X
\end{aligned}
\right.
\]
% foot: Version 9 septembre 2001

\printedpage{18}
% head: A. Grothendieck : Dérivateurs, Ch. V, Édition des Universités de Montpellier II et Paris VII
\noindent\msp{page 28}

\section*{3\quad Passage to categories of diagrams}

We should have denoted the category $X$ by $\mathcal{M}$, for later on it will play the role of a ``category of models''. On the other hand, we shall denote by $X$ the ``index categories'' on which we want to construct the ``categories of diagrams'' of type $X$
\[
\mathcal{M}(X) = \uHom(X,\mathcal{M}).
\]
We thus take
\[
\underbrace{\Sigma}_{\text{quasi-isomorphisms}} \subseteq \mathrm{Fl}(\mathcal{M}), \qquad \Sigma \supset \text{set of identity arrows},
\]
and for every category $X$ we define
\[
\Sigma_X \subset \mathrm{Fl}\,\mathcal{M}(X)
\]
as follows: if $F, G \in \mathrm{Ob}\,\mathcal{M}(X)$ ($F, G : X \rightrightarrows \mathcal{M}$), an arrow $u : F \longrightarrow G$ (homomorphism of functors) is in $\Sigma_X$ if and only if $\forall x \in X$, $u(x) : F(x) \longrightarrow G(x)$ is in $\Sigma$ (``criterion for quasi-isomorphy fiber by fiber'', for the homomorphism between the copresheaves $F$, $G$ under consideration). % [sic] homomophisme -> homomorphisme, letter r missing

Let
\[
a, b \in \mathrm{Ob}\,\mathcal{M}
\]
and consider the corresponding constant functors on $X$
\[
a_X, b_X \in \mathrm{Ob}\,\mathcal{M}(X).
\]

\msp{page 29}We propose to apply the theory of \S2 to the ``calculation'' of
\[
\Hom_{\Sigma_X}(a_X, b_X),
\]
we thus have a canonical bijection
\[
\boxed{\Hom_{\Sigma_X}(a_X, b_X) \simeq \pi_0\uCh_{\Sigma_X}(\mathcal{M}(X); a_x, b_x).} % [sic] a_x, b_x -> a_X, b_X, lowercase index against a_X, b_X elsewhere
\]
We propose to make explicit the category of paths on the right-hand side. We have, leaving aside the data of $\Sigma$, $a$, $b$ in $\mathcal{M}$, the expression for the category of \textit{all} paths
% foot: Version 9 septembre 2001

\end{document}
